\documentclass[11pt,letterpaper]{article}
\usepackage[margin=1in]{geometry}
\usepackage{amsmath,amssymb,amsthm,mathtools}
\usepackage[T1]{fontenc}
\usepackage{lmodern}
\usepackage[hidelinks]{hyperref}
\newtheorem{theorem}{Theorem}
\newtheorem{proposition}[theorem]{Proposition}
\newcommand{\ii}{\mathrm{i}}
\newcommand{\conv}{\operatorname{conv}}
\newcommand{\dist}{\operatorname{dist}}
\title{A Sharp Bound for Centers of Composite Planar Rotations}
\author{Research audit}
\date{12 October 2026}
\begin{document}
\maketitle
\begin{abstract}
An exact elementary counterexample shows that the center of a composition of
positive planar rotations need not belong to the convex hull of the individual
centers, even when their angle sum is at most $\pi$. We prove the sharp
bound $(D/2)\tan(\Theta/4)$ for the distance of the composite center from that
convex hull. The upper bound and sharpness are fully checked in Lean 4.32.2
with mathlib v4.32.2. Originality remains unestablished. An unfolding application
still needs a quantitative separation margin; this result neither resolves
D\"urer's conjecture nor disproves a prismatoid unfolding theorem.
\end{abstract}

\section{Source and scope}
Proposition 9 of O'Rourke's \emph{Prismatoid Band-Unfolding Revisited},
arXiv:2603.09813v3, states exact convex-hull containment for a composite
rotation center when the total angle is at most $\pi$. Its Lemma 8 uses this
ingredient for chain separation. The inspected HTML and PDF page 14 agree on
the unrestricted wording. The latest listed version is dated 18 September
2026; no later correction was located in the bounded search. The author's
separate hosted copy could not be retrieved reliably.

The author's earlier 2017 addendum explicitly distinguishes a small-angle
weighted-center approximation from the true center, which can lie outside
the hull. The exact rotation formulas themselves are classical; see, for
example, Needham's \emph{Visual Complex Analysis}, pp.~20--21. A failed
ingredient does not refute the later main theorem: a valid chain-specific
replacement may still prove the needed conclusion. This observation concerns only the stated
planar-rotation ingredient.\footnote{Sources:
\url{https://arxiv.org/abs/2603.09813v3} and
\url{https://arxiv.org/abs/1709.02433}.}

\section{Exact counterexamples}
Use positive angles counterclockwise, and write
\[
 R_{c,\theta}(z)=e^{\ii\theta}z+(1-e^{\ii\theta})c.
\]
Take $c_1=(0,0)$, $c_2=(1,0)$ and both angles $\pi/2$. Then
\[
 R_{c_1,\pi/2}(x,y)=(-y,x),\qquad
 R_{c_2,\pi/2}(u,v)=(1-v,u-1),
\]
and hence
\[
 F(x,y)=R_{c_2,\pi/2}R_{c_1,\pi/2}(x,y)=(1-x,-1-y).
\]
The equations $F(x,y)=(x,y)$ force $(x,y)=(1/2,-1/2)$. Every rotation
fixes its center. Thus $F$ cannot be a rotation about any point of the segment
$[c_1,c_2]$. Both input angles are positive and their sum is exactly $\pi$.
Reversing composition order or reversing both angle signs reflects the center
to $(1/2,1/2)$, still outside the segment.

The endpoint $\pi$ is not essential. Choose a first-quadrant angle $\theta$
with $\cos\theta=3/5$ and $\sin\theta=4/5$. Then $2\theta<\pi$, and the
same two-center composition has unique center $(1/2,-1/4)$. More generally,
two equal angles $0<\theta<\pi$ give
\[
 c=\frac{c_1+c_2}{2}
   -\frac12\tan(\theta/2)\,\ii(c_2-c_1).
\]
Thus exact hull containment fails even at arbitrarily small positive angles.

\paragraph{Formal verification boundary.}
The accompanying \texttt{RotationAudit.lean} defines ordinary real-plane
rotations and proves the quarter-turn formula, angle conditions, unique fixed
point, and impossibility of any rotation representation with center on the
segment. It also proves the algebraic fixed-point statement for the rational
matrix example. The reproducible build result is recorded in the separate check
log. The rational
example's strict angle interpretation is not currently part of that Lean
certificate. The quarter-turn certificate successfully elaborated with pinned
Lean 4.32.2 and mathlib v4.32.2; its axiom audit reports only
\texttt{propext}, \texttt{Classical.choice} and \texttt{Quot.sound}.
No polyhedral unfolding has been formalized.

\section{The sharp diameter estimate}
\begin{theorem}[Sharp diameter estimate]
Let $m\ge1$ be finite, let $c_1,\dots,c_m\in\mathbb C$, let $\omega_j\ge0$
be angles in radians, and suppose
$0<\Theta=\sum_j\omega_j\le\pi$. Set
\[
 F=R_{c_m,\omega_m}\circ\cdots\circ R_{c_1,\omega_1},\quad
 C=\conv\{c_1,\dots,c_m\},\quad D=\max_{j,k}|c_j-c_k|.
\]
The unique center $c$ of $F$ satisfies
\[
 \dist(c,C)\le\frac D2\tan(\Theta/4).
\]
For each $\Theta\in(0,\pi]$, two equal-angle rotations about the endpoints
of a segment attain equality.
\end{theorem}

\begin{proof}
Let $t_0=0$, $t_j=\sum_{k\le j}\omega_k$, and
$s_j=\Theta/2-t_j$. Direct composition gives
\[
 c=\sum_{j=1}^m\lambda_jc_j,\qquad
 \lambda_j=\frac{e^{\ii(\Theta-t_j)}(1-e^{\ii\omega_j})}
                     {1-e^{\ii\Theta}}.
\]
The denominator is nonzero. With $\lambda_j=p_j+\ii q_j$,
\begin{align*}
 p_j&=\frac{\sin s_{j-1}-\sin s_j}{2\sin(\Theta/2)},\\
 q_j&=\frac{\cos s_j-\cos s_{j-1}}{2\sin(\Theta/2)}.
\end{align*}
Since $s_j$ decreases within $[-\pi/2,\pi/2]$, every $p_j\ge0$.
Telescoping yields $\sum p_j=1$ and $\sum q_j=0$. Therefore
$w=\sum p_jc_j$ lies in $C$.

The sequence $\cos s_j$ increases and then decreases, including the possible
single step that straddles zero. Its total positive increment is its maximum
minus its initial value. Thus, with $Q=\sum_j\max(q_j,0)$,
\[
 Q\le\frac{1-\cos(\Theta/2)}{2\sin(\Theta/2)}
     =\frac12\tan(\Theta/4).
\]
If $Q=0$, then $c=w$. Otherwise the positive and negative coefficients define
two convex combinations $u,v\in C$ such that
$\sum q_jc_j=Q(u-v)$. Any two convex combinations of a set of diameter $D$
are at distance at most $D$, by applying the triangle inequality to the
double convex combination of their differences. Consequently
\[
 |c-w|=\left|\ii\sum q_jc_j\right|
       =Q|u-v|\le QD\le\frac D2\tan(\Theta/4).
\]
This proves the estimate. For equality, take $c_1=0$, $c_2=D$ and
$\omega_1=\omega_2=\Theta/2$. The center is
$D/2-\ii(D/2)\tan(\Theta/4)$, whose distance from $[0,D]$ is the claimed
bound.
\end{proof}

\paragraph{A usable margin condition.}
If a unit vector $a$ satisfies $a\cdot c_j\le b-\delta$ for all $j$, then
$a\cdot c<b$ whenever $\delta>(D/2)\tan(\Theta/4)$. Apply the estimate
to the displayed convex combination $w$. An unfolding application must still
prove this margin from its actual chain geometry.

\section{Formal theorem and verification scope}
The main theorem in \texttt{RotationCenterBound.lean}
uses a finite list of actual rotations. It
assumes nonnegative individual angles, $0<\Theta\le\pi$, $D\ge0$, and
pairwise center distances at most $D$. It proves that the composite is a
rotation with a unique fixed point $c$, and constructs a hull point $w$ with
$|c-w|\le(D/2)\tan(\Theta/4)$. No coefficient estimate or trigonometric
claim remains as a hypothesis. Choosing the maximum pairwise distance gives
the theorem above.

\texttt{RotationSharpness.lean} proves the matching lower bound for every
point in the two-center hull and existence of an attaining hull point, for
every permitted $\Theta$ and $D\ge0$. The endpoint distance is certified as
exactly $D$. The semantic fixture \texttt{RotationExpectedStatement.lean}
rewrites the same composite using the literal map
$e^{\ii\theta}(z-c)+c$ and the ordinary list sum of angles.

The six-module closure uses only \texttt{propext},
\texttt{Classical.choice}, and \texttt{Quot.sound}. The accompanying
source manifest and bounded replay script pin the checked sources and
compiler assumptions. The formal proof bounds the positive cosine variation
by a clipped monotone potential, avoiding an unformalized unimodality claim.
The preceding half-plane consequence is a mathematical corollary, not an
additional Lean declaration in this package.

\section{Priority and remaining questions}
The proof has also passed a separate mathematical audit. Bounded primary-source
searches did not locate this exact sharp diameter theorem; that does not
establish novelty. The 2017 addendum discusses asymptotic approximation, but
this note does not endorse its printed quantitative formula. No novelty is
claimed for the classical composition formula or elementary counterexample.

D\"urer's general question remains the existence of a successful spanning-tree
cut using actual edges in the original metric of every convex three-polytope.
The present estimate supplies no chain-specific separation margin. An
overlapping tree, an affine deformation, or the correction of a subsidiary
lemma does not decide the general question.
\end{document}
